\documentclass[11pt]{article}

\usepackage{amsmath}
\usepackage{amssymb}
\usepackage{graphicx}
\usepackage{natbib}
\usepackage{hyperref}

\graphicspath{{figures/}{../results/figures/}}

\newcommand{\E}{\mathbb{E}}
\newcommand{\R}{\mathbb{R}}
\newcommand{\ip}[2]{\langle #1, #2 \rangle}
\newcommand{\norm}[1]{\lVert #1 \rVert}

\title{Distribution-Free Sequential Learning for Heavy-Tailed Financial Streams}
\author{Rafli\\ \texttt{rafli@pyhron.com}}
\date{July 2026}

\begin{document}

\maketitle

\begin{abstract}
We study online linear regression on data streams whose noise is heavy-tailed and
possibly adversarially contaminated, assuming only a finite second moment. Classical
online gradient descent is fragile in this regime: a single extreme gradient can wipe out
the progress of thousands of rounds. We propose \emph{RobustOMD}, an online mirror
descent scheme whose gradient-clipping threshold is set adaptively by a distribution-free
robust mean estimator (Catoni's M-estimator, median-of-means, or the trimmed mean)
applied to a sliding window of past gradient norms. We state deviation guarantees for the
underlying estimators, sketch a sublinear regret bound for clipped online gradient
descent under a bounded-second-moment assumption, and evaluate the methods on synthetic
heavy-tailed streams and on the Jane Street Real-Time Market Data Forecasting dataset
(47{,}127{,}338 observations, 79 features, 39 symbols over 1{,}699 trading days).
\end{abstract}

\section{Introduction}

Financial market data are a canonical source of heavy tails: returns and microstructure
responses exhibit power-law behavior, occasional data errors act as adversarial
contamination, and distributional assumptions such as sub-Gaussianity are untenable.
Sequential prediction on such streams therefore calls for \emph{distribution-free}
methods, whose guarantees hold under weak moment assumptions only. A rich line of work
in robust statistics shows that the empirical mean is a poor estimator under heavy tails
and that estimators such as Catoni's M-estimator, median-of-means, and the trimmed mean
achieve sub-Gaussian deviation bounds while assuming only finite variance
\citep{catoni2012challenging,devroye2016sub,lugosi2019mean}. In parallel, work on
stochastic optimization has established that gradient clipping is the key algorithmic
ingredient for heavy-tailed gradient noise
\citep{zhang2020adaptive,cutkosky2021high}.

This paper connects the two threads in the online learning setting. We treat the stream
of gradient norms itself as heavy-tailed data, estimate its central tendency with a
robust mean estimator over a sliding window, and use a multiple of that estimate as an
adaptive clipping threshold inside online mirror descent. The resulting algorithm,
RobustOMD, requires no knowledge of the noise scale, adapts to regime changes in the
stream, and degrades gracefully under contamination. We release an open-source
implementation, \texttt{dfsl}, together with reproducible experiments on synthetic
heavy-tailed streams and the Jane Street Kaggle dataset
\citep{janestreet2024kaggle}.

\section{Related Work}

Online convex optimization originates with the online gradient descent (OGD) analysis of
\citet{zinkevich2003online}, which gives $O(\sqrt{T})$ regret for convex losses with
bounded gradients; comprehensive treatments appear in
\citet{cesabianchi2006prediction} and \citet{hazan2016introduction}. Mirror descent,
which OGD instantiates with the Euclidean mirror map, goes back to
\citet{nemirovski1983problem}.

Robust mean estimation under heavy tails is surveyed by \citet{lugosi2019mean}.
\citet{catoni2012challenging} introduced a soft-truncation M-estimator whose deviations
match the Gaussian case up to constants, assuming only finite variance;
\citet{devroye2016sub} characterize sub-Gaussian mean estimators, including
median-of-means and trimmed-mean constructions. In sequential decision making,
\citet{bubeck2013bandits} showed that plugging robust estimators into bandit algorithms
recovers logarithmic regret under heavy-tailed rewards. For stochastic gradient methods,
\citet{zhang2020adaptive} identified heavy-tailed gradient noise in practice and analyzed
clipped SGD, while \citet{cutkosky2021high} obtained high-probability convergence under
heavy tails via clipping and momentum. Our contribution is to make the clipping
threshold itself a robust statistic of the observed gradient-norm stream, in a fully
online, distribution-free manner.

\section{Problem Setting}

\paragraph{Protocol.} Learning proceeds in rounds $t = 1, \dots, T$. At round $t$ the
learner holds a weight vector $w_t \in \R^d$, observes features $x_t \in \R^d$, predicts
$\hat{y}_t = \ip{w_t}{x_t}$, and then observes the target $y_t \in \R$ and a sample
weight $\omega_t > 0$. The learner suffers the \emph{weighted squared loss}
\begin{equation}
\ell_t(w) \;=\; \omega_t \bigl( \ip{w}{x_t} - y_t \bigr)^2 ,
\end{equation}
evaluated prequentially at $w_t$, i.e., before any update that uses $(x_t, y_t)$.

\paragraph{Regret.} Performance is measured by regret against the best fixed linear
predictor in hindsight,
\begin{equation}
R_T \;=\; \sum_{t=1}^{T} \ell_t(w_t) \;-\; \min_{w \in \mathcal{W}} \sum_{t=1}^{T} \ell_t(w),
\end{equation}
where $\mathcal{W} = \{ w : \norm{w}_2 \le R \}$ is a Euclidean ball. An algorithm
\emph{learns} whenever $R_T = o(T)$; empirically we report the slope of
$\log R_T$ versus $\log T$, with slope $< 1$ certifying sublinearity.

\paragraph{Heavy-tailed noise.} We model $y_t = \ip{w^\star}{x_t} + \varepsilon_t$ and
assume only that the noise has a \emph{finite second moment},
$\E[\varepsilon_t^2] \le \sigma^2 < \infty$. No sub-Gaussianity, boundedness, or symmetry
is assumed; $\varepsilon_t$ may be Student-$t$ with barely more than two degrees of
freedom, or two-sided Pareto. In addition, an $\eta$-fraction of the targets may be
arbitrarily contaminated (shifted by outliers of magnitude $\gg \sigma$). Under these
assumptions the stochastic gradients $g_t = 2 \omega_t (\ip{w_t}{x_t} - y_t) x_t$ are
themselves heavy-tailed, which is precisely the regime in which unclipped OGD fails.

\section{Robust Estimators}

All three estimators below take a sample $z_1, \dots, z_n$ with mean $\mu$ and variance
$\sigma^2 < \infty$ and return an estimate $\hat{\mu}$ with sub-Gaussian-type deviations,
despite the possibly heavy-tailed sample.

\paragraph{Catoni's M-estimator.} With the influence function
\begin{equation}
\psi(t) \;=\; \operatorname{sign}(t) \, \log\!\Bigl( 1 + |t| + \tfrac{t^2}{2} \Bigr),
\end{equation}
Catoni's estimator $\hat{\mu}_C$ is the root in $\theta$ of
$\sum_{i=1}^{n} \psi\bigl( \alpha (z_i - \theta) \bigr) = 0$, a strictly decreasing
function of $\theta$, solvable by bisection. With
$\alpha = \sqrt{2 \log(1/\delta) / (n \sigma^2)}$,
\citet{catoni2012challenging} shows that for $n \gtrsim \log(1/\delta)$,
\begin{equation}
|\hat{\mu}_C - \mu| \;\le\; C \, \sigma \sqrt{\frac{\log(1/\delta)}{n}}
\qquad \text{with probability } \ge 1 - 2\delta,
\end{equation}
matching the deviations of the empirical mean of a Gaussian sample up to constants.

\paragraph{Median-of-means.} Partition the sample into
$k = \lceil 8 \log(1/\delta) \rceil$ blocks of (nearly) equal size, compute each block
mean, and return their median $\hat{\mu}_{\mathrm{MoM}}$. Then
\citep{nemirovski1983problem,devroye2016sub,lugosi2019mean}
\begin{equation}
|\hat{\mu}_{\mathrm{MoM}} - \mu| \;\le\; C \, \sigma \sqrt{\frac{\log(1/\delta)}{n}}
\qquad \text{with probability } \ge 1 - \delta .
\end{equation}
Median-of-means additionally tolerates a constant fraction of adversarially corrupted
blocks, giving it a positive breakdown point.

\paragraph{Trimmed mean.} Sort the sample, discard the $\lfloor \gamma n \rfloor$
smallest and largest points ($0 \le \gamma < 1/2$), and average the rest. For a suitable
$\gamma \asymp \log(1/\delta)/n + \eta$ under $\eta$-contamination, the trimmed mean is a
sub-Gaussian estimator that is simultaneously robust to contamination
\citep{lugosi2019mean,devroye2016sub}.

\section{Algorithms}

The baseline is projected OGD \citep{zinkevich2003online} with step size
$\eta_t = \eta / \sqrt{t}$. Our robust variants keep the OGD template and modify only
the gradient through a clipping hook. RobustOMD is online mirror descent with the
Euclidean mirror map $\Phi(w) = \tfrac{1}{2}\norm{w}_2^2$ and a distribution-free
clipping threshold; one round proceeds as follows.

\begin{enumerate}
  \item Observe $x_t$; predict $\hat{y}_t = \ip{w_t}{x_t}$; observe $y_t, \omega_t$ and
        incur $\ell_t = \omega_t (\hat{y}_t - y_t)^2$.
  \item Compute the gradient $g_t = 2 \omega_t (\hat{y}_t - y_t) \, x_t$ and append
        $\norm{g_t}_2$ to a sliding window $H_t$ of the most recent $W$ unclipped
        gradient norms.
  \item Estimate the central gradient magnitude
        $\hat{m}_t = \mathrm{RobustMean}(H_t)$, where $\mathrm{RobustMean}$ is Catoni's
        M-estimator, median-of-means, or the trimmed mean, and set the threshold
        $\tau_t = c \, \hat{m}_t$ for a clip multiplier $c > 0$.
  \item Clip: $\tilde{g}_t = g_t \cdot \min\{1, \tau_t / \norm{g_t}_2\}$.
  \item Mirror-descent step and projection:
        $w_{t+1} = \Pi_{\mathcal{W}} \bigl( w_t - \eta_t \, \tilde{g}_t \bigr)$.
\end{enumerate}

The \emph{AdaptiveClip} baseline replaces step 3 by an empirical quantile of $H_t$
(threshold $\tau_t = \mathrm{Quantile}_q(H_t)$), a natural but less principled choice:
sample quantiles of a heavy-tailed norm sequence have high variance, whereas the robust
mean estimators concentrate at the sub-Gaussian rate regardless of the tail index, as
long as the norms have finite variance. A regret bound for the clipped update under the
bounded-second-moment assumption is sketched in Appendix~\ref{app:proofs}.

\section{Experiments}

\paragraph{Synthetic heavy-tailed streams.} We generate $y_t = \ip{w^\star}{x_t} +
\varepsilon_t$ with $x_t \sim \mathcal{N}(0, I_{10})$, unit-norm $w^\star$, and noise
drawn from Student-$t$ ($\nu = 2.5$, finite variance but infinite higher moments),
two-sided Pareto, or Cauchy distributions; optionally a fraction
$\eta \in \{0, 0.05, 0.1\}$ of targets is shifted by $\pm 50$. We compare OGD,
AdaptiveClip, and RobustOMD with each estimator, all sharing learning rates, and report
cumulative regret against the best fixed ridge solution, rolling prequential loss, and
breakdown curves under increasing contamination. Regret and loss curves are produced by
\texttt{experiments/benchmark.py} and stored under \texttt{results/benchmark}
(\texttt{regret\_curves.png}, \texttt{loss\_curves.png}); ablation heatmaps over the
estimator, window size $W \in \{64, 256, 1024\}$, and clip multiplier
$c \in \{1, 3, 5\}$ are produced by \texttt{experiments/ablation.py}. Across noise
settings, both robust variants track the comparator with visibly sublinear regret
(empirical slope $< 1$), while plain OGD suffers loss spikes aligned with extreme noise
draws; under contamination the gap widens and RobustOMD is the most stable, consistent
with the breakdown properties of its estimators.

\paragraph{Jane Street Real-Time Market Data Forecasting.} The real-data testbed is the
Kaggle competition dataset \citep{janestreet2024kaggle}: 47{,}127{,}338 rows of
anonymized market observations covering 1{,}699 trading days (\texttt{date\_id}
0--1698), 39 symbols, and up to 968 intraday time steps per day, with 79 features
(47 of which contain nulls, up to 21.8\% missing) and 9 responders clipped to
$[-5, 5]$. The competition target is \texttt{responder\_6} and the metric is the
sample-weighted zero-mean coefficient of determination
\begin{equation}
R^2 \;=\; 1 - \frac{\sum_i \omega_i (y_i - \hat{y}_i)^2}{\sum_i \omega_i y_i^2},
\end{equation}
with sample weights $\omega_i > 0$ supplied by the dataset. We stream the data in
chronological order (sorted by \texttt{date\_id}, \texttt{time\_id},
\texttt{symbol\_id}), fill feature nulls with zero, and run the same learner grid;
dataset characterization figures (missingness over time, weight distribution, intraday
seasonality, regime shifts) are generated by the analysis notebooks and stored under
\texttt{results/figures}. The heavy-tailed, weakly predictable responders make this a
demanding distribution-free benchmark: robust clipping mainly prevents rare extreme
rounds from destroying the accumulated fit, which shows up as a higher weighted $R^2$
over long horizons.

\section{Conclusion}

We presented a distribution-free approach to sequential regression on heavy-tailed
streams: use robust mean estimators, whose sub-Gaussian deviation guarantees require
only finite variance, to set the clipping threshold of online mirror descent. The
approach needs no prior knowledge of the noise scale, inherits the breakdown robustness
of its estimator, and is simple to implement — the \texttt{dfsl} library, configs, and
scripts reproduce every experiment in this paper. Future work includes high-probability
regret bounds for the fully adaptive threshold, per-coordinate robust clipping, and
extensions to nonstationary comparators, which the regime-shift structure of financial
streams clearly motivates.

\appendix
% Appendix for "Distribution-Free Sequential Learning for Heavy-Tailed Financial
% Streams". Included from main.tex via \appendix % Appendix for "Distribution-Free Sequential Learning for Heavy-Tailed Financial
% Streams". Included from main.tex via \appendix % Appendix for "Distribution-Free Sequential Learning for Heavy-Tailed Financial
% Streams". Included from main.tex via \appendix \input{appendix}; contains no preamble.

\section{Proof Sketches}
\label{app:proofs}

\subsection{Regret of clipped OGD under a bounded second moment}
\label{app:clipped-ogd}

We sketch why clipping with a well-chosen threshold preserves $O(\sqrt{T})$ regret in
expectation when gradients have only a finite second moment. Let
$\mathcal{W} = \{ w : \norm{w}_2 \le R \}$, let $g_t$ denote the (unclipped) gradient of
the weighted squared loss at $w_t$, and assume $\E \bigl[ \norm{g_t}_2^2 \mid
\mathcal{F}_{t-1} \bigr] \le G^2$ for all $t$. Fix a threshold $\tau > 0$ and let
$\tilde{g}_t = g_t \min\{1, \tau / \norm{g_t}_2\}$ be the clipped gradient used in the
update $w_{t+1} = \Pi_{\mathcal{W}}(w_t - \eta_t \tilde{g}_t)$.

\paragraph{Step 1: descent inequality for the clipped sequence.} The standard OGD
argument \citep{zinkevich2003online,hazan2016introduction} applied to the surrogate
linear losses $w \mapsto \ip{\tilde{g}_t}{w}$ gives, for any comparator
$u \in \mathcal{W}$ and $\eta_t = \eta/\sqrt{t}$,
\begin{equation}
\sum_{t=1}^{T} \ip{\tilde{g}_t}{w_t - u}
\;\le\; \frac{(2R)^2}{2\eta_T} + \sum_{t=1}^{T} \frac{\eta_t}{2} \norm{\tilde{g}_t}_2^2
\;\le\; \frac{2R^2 \sqrt{T}}{\eta} + \eta \tau^2 \sqrt{T},
\end{equation}
because $\norm{\tilde{g}_t}_2 \le \tau$ deterministically.

\paragraph{Step 2: bias of clipping.} By convexity of $\ell_t$,
$\ell_t(w_t) - \ell_t(u) \le \ip{g_t}{w_t - u}$, so it remains to control the clipping
bias $b_t = \ip{g_t - \tilde{g}_t}{w_t - u}$. Since
$\norm{g_t - \tilde{g}_t}_2 = (\norm{g_t}_2 - \tau)_+ \le
\norm{g_t}_2 \mathbf{1}\{\norm{g_t}_2 > \tau\}$ and $\norm{w_t - u}_2 \le 2R$,
Cauchy--Schwarz followed by Markov's inequality yields
\begin{equation}
\E[b_t]
\;\le\; 2R \, \E\bigl[ \norm{g_t}_2 \mathbf{1}\{ \norm{g_t}_2 > \tau \} \bigr]
\;\le\; 2R \, \frac{\E[\norm{g_t}_2^2]}{\tau}
\;\le\; \frac{2 R G^2}{\tau}.
\end{equation}
Only the second moment of the gradient is used; no higher moments and no sub-Gaussian
tail are required.

\paragraph{Step 3: balancing.} Summing over $t$ and combining Steps 1--2,
\begin{equation}
\E[R_T] \;\le\; \frac{2R^2\sqrt{T}}{\eta} + \eta \tau^2 \sqrt{T} + \frac{2RG^2 T}{\tau}.
\end{equation}
Choosing $\tau \asymp G\,T^{1/4}$ and $\eta \asymp R / (G\,T^{1/4})$ balances the three
terms and gives $\E[R_T] = O\!\bigl( R G \, T^{3/4} \bigr)$ under the bare second-moment
assumption; when the gradient norms additionally have bounded variance around a mean of
order $m$, taking $\tau = c\,m$ with a constant multiplier $c$ recovers the familiar
$O(RG\sqrt{T})$ rate up to the (bounded) fraction of clipped rounds. High-probability
versions follow from the truncation-plus-martingale arguments of
\citet{cutkosky2021high} and \citet{zhang2020adaptive}.

\subsection{Validity of the data-driven threshold}
\label{app:threshold}

RobustOMD replaces the oracle threshold $\tau$ by $\tau_t = c\,\hat{m}_t$, where
$\hat{m}_t$ is a robust mean estimate of the last $W$ unclipped gradient norms. If the
norms in the window are (conditionally) i.i.d.\ with mean $m$ and variance
$s^2 < \infty$, then Catoni's estimator and median-of-means both satisfy
\begin{equation}
| \hat{m}_t - m | \;\le\; C\, s \sqrt{\frac{\log(1/\delta)}{W}}
\qquad \text{with probability } \ge 1 - 2\delta,
\end{equation}
\citep{catoni2012challenging,devroye2016sub,lugosi2019mean}, so $\tau_t$ concentrates in
the constant-multiple band $[\,c\,m - \epsilon_W,\; c\,m + \epsilon_W\,]$ with
$\epsilon_W \to 0$ as the window grows. Substituting this band into the analysis of
Appendix~\ref{app:clipped-ogd} changes the bound only by constants. Under
$\eta$-contamination of the stream, median-of-means and the trimmed mean keep this
guarantee as long as the corrupted fraction of the window stays below their breakdown
points, which explains the flat breakdown curves observed empirically; a sample quantile
(AdaptiveClip) enjoys no comparable deviation guarantee under heavy tails, consistent
with its higher variance in the ablations \citep{bubeck2013bandits}.

\section{Additional Experiment Details}
\label{app:experiments}

\paragraph{Synthetic streams.} Unless stated otherwise: horizon $T = 20{,}000$,
dimension $d = 10$, features $x_t \sim \mathcal{N}(0, I_d)$, unit-norm ground truth
$w^\star$, Student-$t$ noise with $\nu = 2.5$ and unit scale, generated from a fixed
numpy \texttt{Generator} seed (0). Contaminated runs shift a uniformly chosen 5\% or
10\% of targets by $\pm 50$. All learners share the learning rate $\eta = 0.1$; robust
variants use window $W = 256$ and clip multiplier $c = 3$; AdaptiveClip uses the 0.9
empirical quantile. Every configuration is fully specified by the YAML files in
\texttt{configs/} and reproduced by \texttt{scripts/reproduce.sh}.

\paragraph{Ablation grid.} The sweep of \texttt{experiments/ablation.py} crosses the
estimator $\in \{\text{catoni}, \text{median-of-means}, \text{trimmed mean}\}$, window
$W \in \{64, 256, 1024\}$, and multiplier $c \in \{1, 3, 5\}$ on the contaminated
Student-$t$ stream ($\eta = 0.1$). Small windows make the threshold noisy, very large
$c$ effectively disables clipping, and $c = 1$ over-clips informative gradients; the
plateau around $(W, c) = (256, 3)$ is broad for all three estimators, indicating that no
fine tuning is needed.

\paragraph{Jane Street preprocessing.} The stream is built from the hive-partitioned
\texttt{train.parquet} via lazy polars scans: filter to the requested \texttt{date\_id}
range, drop rows with a null target, keep the 79 features plus
(\texttt{date\_id}, \texttt{time\_id}, \texttt{symbol\_id}, \texttt{weight}), cast
features to \texttt{Float64}, fill nulls with $0$, and sort chronologically by
(\texttt{date\_id}, \texttt{time\_id}, \texttt{symbol\_id}). The default real-data
config (\texttt{configs/jane.yaml}) uses the first 31 trading days capped at 200{,}000
rows, learning rate $0.05$, window $512$, and $c = 3$; metrics are always computed with
the competition's sample weights. Runs write \texttt{run.parquet},
\texttt{metrics.json}, and the resolved \texttt{config.yaml} for exact reproducibility.
; contains no preamble.

\section{Proof Sketches}
\label{app:proofs}

\subsection{Regret of clipped OGD under a bounded second moment}
\label{app:clipped-ogd}

We sketch why clipping with a well-chosen threshold preserves $O(\sqrt{T})$ regret in
expectation when gradients have only a finite second moment. Let
$\mathcal{W} = \{ w : \norm{w}_2 \le R \}$, let $g_t$ denote the (unclipped) gradient of
the weighted squared loss at $w_t$, and assume $\E \bigl[ \norm{g_t}_2^2 \mid
\mathcal{F}_{t-1} \bigr] \le G^2$ for all $t$. Fix a threshold $\tau > 0$ and let
$\tilde{g}_t = g_t \min\{1, \tau / \norm{g_t}_2\}$ be the clipped gradient used in the
update $w_{t+1} = \Pi_{\mathcal{W}}(w_t - \eta_t \tilde{g}_t)$.

\paragraph{Step 1: descent inequality for the clipped sequence.} The standard OGD
argument \citep{zinkevich2003online,hazan2016introduction} applied to the surrogate
linear losses $w \mapsto \ip{\tilde{g}_t}{w}$ gives, for any comparator
$u \in \mathcal{W}$ and $\eta_t = \eta/\sqrt{t}$,
\begin{equation}
\sum_{t=1}^{T} \ip{\tilde{g}_t}{w_t - u}
\;\le\; \frac{(2R)^2}{2\eta_T} + \sum_{t=1}^{T} \frac{\eta_t}{2} \norm{\tilde{g}_t}_2^2
\;\le\; \frac{2R^2 \sqrt{T}}{\eta} + \eta \tau^2 \sqrt{T},
\end{equation}
because $\norm{\tilde{g}_t}_2 \le \tau$ deterministically.

\paragraph{Step 2: bias of clipping.} By convexity of $\ell_t$,
$\ell_t(w_t) - \ell_t(u) \le \ip{g_t}{w_t - u}$, so it remains to control the clipping
bias $b_t = \ip{g_t - \tilde{g}_t}{w_t - u}$. Since
$\norm{g_t - \tilde{g}_t}_2 = (\norm{g_t}_2 - \tau)_+ \le
\norm{g_t}_2 \mathbf{1}\{\norm{g_t}_2 > \tau\}$ and $\norm{w_t - u}_2 \le 2R$,
Cauchy--Schwarz followed by Markov's inequality yields
\begin{equation}
\E[b_t]
\;\le\; 2R \, \E\bigl[ \norm{g_t}_2 \mathbf{1}\{ \norm{g_t}_2 > \tau \} \bigr]
\;\le\; 2R \, \frac{\E[\norm{g_t}_2^2]}{\tau}
\;\le\; \frac{2 R G^2}{\tau}.
\end{equation}
Only the second moment of the gradient is used; no higher moments and no sub-Gaussian
tail are required.

\paragraph{Step 3: balancing.} Summing over $t$ and combining Steps 1--2,
\begin{equation}
\E[R_T] \;\le\; \frac{2R^2\sqrt{T}}{\eta} + \eta \tau^2 \sqrt{T} + \frac{2RG^2 T}{\tau}.
\end{equation}
Choosing $\tau \asymp G\,T^{1/4}$ and $\eta \asymp R / (G\,T^{1/4})$ balances the three
terms and gives $\E[R_T] = O\!\bigl( R G \, T^{3/4} \bigr)$ under the bare second-moment
assumption; when the gradient norms additionally have bounded variance around a mean of
order $m$, taking $\tau = c\,m$ with a constant multiplier $c$ recovers the familiar
$O(RG\sqrt{T})$ rate up to the (bounded) fraction of clipped rounds. High-probability
versions follow from the truncation-plus-martingale arguments of
\citet{cutkosky2021high} and \citet{zhang2020adaptive}.

\subsection{Validity of the data-driven threshold}
\label{app:threshold}

RobustOMD replaces the oracle threshold $\tau$ by $\tau_t = c\,\hat{m}_t$, where
$\hat{m}_t$ is a robust mean estimate of the last $W$ unclipped gradient norms. If the
norms in the window are (conditionally) i.i.d.\ with mean $m$ and variance
$s^2 < \infty$, then Catoni's estimator and median-of-means both satisfy
\begin{equation}
| \hat{m}_t - m | \;\le\; C\, s \sqrt{\frac{\log(1/\delta)}{W}}
\qquad \text{with probability } \ge 1 - 2\delta,
\end{equation}
\citep{catoni2012challenging,devroye2016sub,lugosi2019mean}, so $\tau_t$ concentrates in
the constant-multiple band $[\,c\,m - \epsilon_W,\; c\,m + \epsilon_W\,]$ with
$\epsilon_W \to 0$ as the window grows. Substituting this band into the analysis of
Appendix~\ref{app:clipped-ogd} changes the bound only by constants. Under
$\eta$-contamination of the stream, median-of-means and the trimmed mean keep this
guarantee as long as the corrupted fraction of the window stays below their breakdown
points, which explains the flat breakdown curves observed empirically; a sample quantile
(AdaptiveClip) enjoys no comparable deviation guarantee under heavy tails, consistent
with its higher variance in the ablations \citep{bubeck2013bandits}.

\section{Additional Experiment Details}
\label{app:experiments}

\paragraph{Synthetic streams.} Unless stated otherwise: horizon $T = 20{,}000$,
dimension $d = 10$, features $x_t \sim \mathcal{N}(0, I_d)$, unit-norm ground truth
$w^\star$, Student-$t$ noise with $\nu = 2.5$ and unit scale, generated from a fixed
numpy \texttt{Generator} seed (0). Contaminated runs shift a uniformly chosen 5\% or
10\% of targets by $\pm 50$. All learners share the learning rate $\eta = 0.1$; robust
variants use window $W = 256$ and clip multiplier $c = 3$; AdaptiveClip uses the 0.9
empirical quantile. Every configuration is fully specified by the YAML files in
\texttt{configs/} and reproduced by \texttt{scripts/reproduce.sh}.

\paragraph{Ablation grid.} The sweep of \texttt{experiments/ablation.py} crosses the
estimator $\in \{\text{catoni}, \text{median-of-means}, \text{trimmed mean}\}$, window
$W \in \{64, 256, 1024\}$, and multiplier $c \in \{1, 3, 5\}$ on the contaminated
Student-$t$ stream ($\eta = 0.1$). Small windows make the threshold noisy, very large
$c$ effectively disables clipping, and $c = 1$ over-clips informative gradients; the
plateau around $(W, c) = (256, 3)$ is broad for all three estimators, indicating that no
fine tuning is needed.

\paragraph{Jane Street preprocessing.} The stream is built from the hive-partitioned
\texttt{train.parquet} via lazy polars scans: filter to the requested \texttt{date\_id}
range, drop rows with a null target, keep the 79 features plus
(\texttt{date\_id}, \texttt{time\_id}, \texttt{symbol\_id}, \texttt{weight}), cast
features to \texttt{Float64}, fill nulls with $0$, and sort chronologically by
(\texttt{date\_id}, \texttt{time\_id}, \texttt{symbol\_id}). The default real-data
config (\texttt{configs/jane.yaml}) uses the first 31 trading days capped at 200{,}000
rows, learning rate $0.05$, window $512$, and $c = 3$; metrics are always computed with
the competition's sample weights. Runs write \texttt{run.parquet},
\texttt{metrics.json}, and the resolved \texttt{config.yaml} for exact reproducibility.
; contains no preamble.

\section{Proof Sketches}
\label{app:proofs}

\subsection{Regret of clipped OGD under a bounded second moment}
\label{app:clipped-ogd}

We sketch why clipping with a well-chosen threshold preserves $O(\sqrt{T})$ regret in
expectation when gradients have only a finite second moment. Let
$\mathcal{W} = \{ w : \norm{w}_2 \le R \}$, let $g_t$ denote the (unclipped) gradient of
the weighted squared loss at $w_t$, and assume $\E \bigl[ \norm{g_t}_2^2 \mid
\mathcal{F}_{t-1} \bigr] \le G^2$ for all $t$. Fix a threshold $\tau > 0$ and let
$\tilde{g}_t = g_t \min\{1, \tau / \norm{g_t}_2\}$ be the clipped gradient used in the
update $w_{t+1} = \Pi_{\mathcal{W}}(w_t - \eta_t \tilde{g}_t)$.

\paragraph{Step 1: descent inequality for the clipped sequence.} The standard OGD
argument \citep{zinkevich2003online,hazan2016introduction} applied to the surrogate
linear losses $w \mapsto \ip{\tilde{g}_t}{w}$ gives, for any comparator
$u \in \mathcal{W}$ and $\eta_t = \eta/\sqrt{t}$,
\begin{equation}
\sum_{t=1}^{T} \ip{\tilde{g}_t}{w_t - u}
\;\le\; \frac{(2R)^2}{2\eta_T} + \sum_{t=1}^{T} \frac{\eta_t}{2} \norm{\tilde{g}_t}_2^2
\;\le\; \frac{2R^2 \sqrt{T}}{\eta} + \eta \tau^2 \sqrt{T},
\end{equation}
because $\norm{\tilde{g}_t}_2 \le \tau$ deterministically.

\paragraph{Step 2: bias of clipping.} By convexity of $\ell_t$,
$\ell_t(w_t) - \ell_t(u) \le \ip{g_t}{w_t - u}$, so it remains to control the clipping
bias $b_t = \ip{g_t - \tilde{g}_t}{w_t - u}$. Since
$\norm{g_t - \tilde{g}_t}_2 = (\norm{g_t}_2 - \tau)_+ \le
\norm{g_t}_2 \mathbf{1}\{\norm{g_t}_2 > \tau\}$ and $\norm{w_t - u}_2 \le 2R$,
Cauchy--Schwarz followed by Markov's inequality yields
\begin{equation}
\E[b_t]
\;\le\; 2R \, \E\bigl[ \norm{g_t}_2 \mathbf{1}\{ \norm{g_t}_2 > \tau \} \bigr]
\;\le\; 2R \, \frac{\E[\norm{g_t}_2^2]}{\tau}
\;\le\; \frac{2 R G^2}{\tau}.
\end{equation}
Only the second moment of the gradient is used; no higher moments and no sub-Gaussian
tail are required.

\paragraph{Step 3: balancing.} Summing over $t$ and combining Steps 1--2,
\begin{equation}
\E[R_T] \;\le\; \frac{2R^2\sqrt{T}}{\eta} + \eta \tau^2 \sqrt{T} + \frac{2RG^2 T}{\tau}.
\end{equation}
Choosing $\tau \asymp G\,T^{1/4}$ and $\eta \asymp R / (G\,T^{1/4})$ balances the three
terms and gives $\E[R_T] = O\!\bigl( R G \, T^{3/4} \bigr)$ under the bare second-moment
assumption; when the gradient norms additionally have bounded variance around a mean of
order $m$, taking $\tau = c\,m$ with a constant multiplier $c$ recovers the familiar
$O(RG\sqrt{T})$ rate up to the (bounded) fraction of clipped rounds. High-probability
versions follow from the truncation-plus-martingale arguments of
\citet{cutkosky2021high} and \citet{zhang2020adaptive}.

\subsection{Validity of the data-driven threshold}
\label{app:threshold}

RobustOMD replaces the oracle threshold $\tau$ by $\tau_t = c\,\hat{m}_t$, where
$\hat{m}_t$ is a robust mean estimate of the last $W$ unclipped gradient norms. If the
norms in the window are (conditionally) i.i.d.\ with mean $m$ and variance
$s^2 < \infty$, then Catoni's estimator and median-of-means both satisfy
\begin{equation}
| \hat{m}_t - m | \;\le\; C\, s \sqrt{\frac{\log(1/\delta)}{W}}
\qquad \text{with probability } \ge 1 - 2\delta,
\end{equation}
\citep{catoni2012challenging,devroye2016sub,lugosi2019mean}, so $\tau_t$ concentrates in
the constant-multiple band $[\,c\,m - \epsilon_W,\; c\,m + \epsilon_W\,]$ with
$\epsilon_W \to 0$ as the window grows. Substituting this band into the analysis of
Appendix~\ref{app:clipped-ogd} changes the bound only by constants. Under
$\eta$-contamination of the stream, median-of-means and the trimmed mean keep this
guarantee as long as the corrupted fraction of the window stays below their breakdown
points, which explains the flat breakdown curves observed empirically; a sample quantile
(AdaptiveClip) enjoys no comparable deviation guarantee under heavy tails, consistent
with its higher variance in the ablations \citep{bubeck2013bandits}.

\section{Additional Experiment Details}
\label{app:experiments}

\paragraph{Synthetic streams.} Unless stated otherwise: horizon $T = 20{,}000$,
dimension $d = 10$, features $x_t \sim \mathcal{N}(0, I_d)$, unit-norm ground truth
$w^\star$, Student-$t$ noise with $\nu = 2.5$ and unit scale, generated from a fixed
numpy \texttt{Generator} seed (0). Contaminated runs shift a uniformly chosen 5\% or
10\% of targets by $\pm 50$. All learners share the learning rate $\eta = 0.1$; robust
variants use window $W = 256$ and clip multiplier $c = 3$; AdaptiveClip uses the 0.9
empirical quantile. Every configuration is fully specified by the YAML files in
\texttt{configs/} and reproduced by \texttt{scripts/reproduce.sh}.

\paragraph{Ablation grid.} The sweep of \texttt{experiments/ablation.py} crosses the
estimator $\in \{\text{catoni}, \text{median-of-means}, \text{trimmed mean}\}$, window
$W \in \{64, 256, 1024\}$, and multiplier $c \in \{1, 3, 5\}$ on the contaminated
Student-$t$ stream ($\eta = 0.1$). Small windows make the threshold noisy, very large
$c$ effectively disables clipping, and $c = 1$ over-clips informative gradients; the
plateau around $(W, c) = (256, 3)$ is broad for all three estimators, indicating that no
fine tuning is needed.

\paragraph{Jane Street preprocessing.} The stream is built from the hive-partitioned
\texttt{train.parquet} via lazy polars scans: filter to the requested \texttt{date\_id}
range, drop rows with a null target, keep the 79 features plus
(\texttt{date\_id}, \texttt{time\_id}, \texttt{symbol\_id}, \texttt{weight}), cast
features to \texttt{Float64}, fill nulls with $0$, and sort chronologically by
(\texttt{date\_id}, \texttt{time\_id}, \texttt{symbol\_id}). The default real-data
config (\texttt{configs/jane.yaml}) uses the first 31 trading days capped at 200{,}000
rows, learning rate $0.05$, window $512$, and $c = 3$; metrics are always computed with
the competition's sample weights. Runs write \texttt{run.parquet},
\texttt{metrics.json}, and the resolved \texttt{config.yaml} for exact reproducibility.


\bibliographystyle{plainnat}
\bibliography{references}

\end{document}
